\documentclass[10pt, ukrainian]{article}
\setlength\textwidth{150mm} \setlength\textheight{250mm}
\pagestyle{empty}
\usepackage[T2A]{fontenc}
\usepackage[utf8]{inputenc}
\usepackage[ukrainian]{babel}
\usepackage{graphicx}
\usepackage{caption}
\DeclareCaptionFormat{nocaption}{}
\captionsetup{format=nocaption,aboveskip=0pt,belowskip=0pt}
\usepackage[Export]{adjustbox}
\adjustboxset{max size={0.9\linewidth}{0.9\paperheight}}
\begin{document}
{
\begin {Large}

1. Написати функцію, що перевіряє чи є послідовність дійсних чисел монотонною.

2. Розглядаються неспадні послідовності цілих чисел, що зберігаються в списках. Написати функцію, що утворює нову неспадну послідовність, яка є злиттям двох своїх аргументів. Якщо елемент в першу послідовність входить $k_1$ раз, а в другу --- $k_2$ разів, то в результат він має входити $k_1+k_2$ разів. Наприклад, за послідовностей 1, 3, 3, 5 та 2, 3, 5, 8 результатом має бути 1, 2, 3, 3, 3, 5, 5, 8.

3. Розглядаються неспадні послідовності цілих чисел, що зберігаються в списках. Написати функцію, що утворює нову неспадну послідовність, яка є об'єднанням двох своїх аргументів. Якщо елемент в першу послідовність входить $k_1$ раз, а в другу --- $k_2$ разів, то в результат він має входити $max\left\{k_1,k_2\right\}$ разів. Наприклад, за послідовностей 1, 3, 3, 5 та 2, 3, 5, 8 результатом має бути 1, 2, 3, 3, 5, 8.

4. Розглядаються неспадні послідовності цілих чисел, що зберігаються в списках. Написати функцію, що утворює нову неспадну послідовність, яка є перетином двох своїх аргументів. Якщо елемент в першу послідовність входить $k_1$ раз, а в другу --- $k_2$ разів, то в результат він має входити $min\left\{k_1,k_2\right\}$ разів. Наприклад, за послідовностей 1, 3, 3, 5 та 2, 3, 3, 3, 5, 8 результатом має бути 3, 3, 5.

5. Розглядаються неспадні послідовності цілих чисел, що зберігаються в списках. Написати функцію, що перевіряє чи складаються дві послідовності з однакових елементів. Кількість входжень елемента в послідовність до уваги не брати. Наприклад, послідовності 1, 1, 3 та 1, 1, 1, 3, 3 складаються з однакових елементів. Обмеження: використовувати додаткові структури даних для збереження послідовностей або їх частин заборонено.

6. Список містить цілі числа. Вилучити зі списку усі повторення (кожне значення має залишитися в одному екземплярі). Використання додаткових структур даних заборонено. Для кожного елемента в списку має залишитись його перше входження і при цьому порядок елементів змінювати не можна.

\newpage

7. Список містить цілі числа. Вилучити зі списку усі повторення (кожне значення має залишитися в одному екземплярі). Використання додаткових структур даних заборонено. Для кожного елемента в списку має залишитись його останнє входження і при цьому порядок елементів змінювати не можна.

8. Рядок містить слова, відокремлені крапками та комами. Знайти довжину найдовшого слова рядка. Копій частини рядка довших 1 символу створювати не можна.

9. Рядок містить слова, відокремлені крапками та комами. Знайти останнє найдовше слово рядка. Копії частини рядка довші 1 символу можна створювати тільки в інструкції \textbf{return}.

10. Рядок містить слова, відокремлені крапками та комами. Знайти кількість слів, довжина яких більша 3. Копій частини рядка довших 1 символу створювати не можна.

11. Рядок містить слова, відокремлені крапками та комами. Знайти кількість слів, довжина яких дорівнює 13. Копій частини рядка довших 1 символу створювати не можна.

12. Рядок містить слова, відокремлені крапками та комами. Знайти кількість слів, що містять кожну з літер s, e, l, f. Копій частини рядка довших 1 символу створювати не можна.

13. Рядок містить слова, відокремлені крапками та комами. Знайти останнє слово довжини 13. Копії частини рядка довші 1 символу можна створювати тільки в інструкції \textbf{return}.

14. Написати функцію, що стискає послідовність цілих, зображену списком: нулі в кінці та на початку вилучаються, у середині послідовності кожна послідовність нулів замінюється на один нуль. Наприклад, з послідовності 0,0,1,2,0,0,0,0,3,0,5,0,0 має стати 1,2,0,3,0,5. Нові об'єкти, що зберігають послідовності утворювати не можна.

15. Бінарні відношення $R_1$, $R_2$ задано на множині $\overline{0,n-1}$ матрицями відношення з цілими елементами. Перевірити, чи правда, що  $R_1\subseteq R_2$. Обов'язково зазначити, яке зображення матриці використовується в розв'язанні: двовимірний масив типу \textbf{ndarray} або список списків.

\newpage

16. Бінарні відношення $R_1$, $R_2$ задано на множині $\overline{0,n-1}$ матрицями відношення з цілими елементами. Перевірити, чи правда, що $R_1\subset R_2$. Обов'язково зазначити, яке зображення матриці використовується в розв'язанні: двовимірний масив типу \textbf{ndarray} або список списків.

17. Бінарні відношення $R_1$, $R_2$ задано на множині $\overline{0,n-1}$ матрицями відношення з цілими елементами. Обчислити композицію $R_1\circ R_2$ відношень. Обов'язково зазначити, яке зображення матриці використовується в розв'язанні: двовимірний масив типу \textbf{ndarray} або список списків.

18. Бінарні відношення $R_1$, $R_2$ задано на множині $\overline{0,n-1}$ матрицями відношення з цілими елементами. Обчислити матрицю відношення $R_1\cup R_2$. Обов'язково зазначити, яке зображення матриці використовується в розв'язанні: двовимірний масив типу \textbf{ndarray} або список списків. Для всіх матриць має використовуватись однакове зображення.

19. Бінарні відношення $R_1$, $R_2$ задано на множині $\overline{0,n-1}$ матрицями відношення з цілими елементами. Обчислити матрицю відношення $R_1\cap R_2$. Обов'язково зазначити, яке зображення матриці використовується в розв'язанні: двовимірний масив типу \textbf{ndarray} або список списків. Для всіх матриць має використовуватись однакове зображення.

20. Бінарні відношення $R_1$, $R_2$ задано на множині $\overline{0,n-1}$ матрицями відношення з цілими елементами. Обчислити матрицю відношення $R_1\setminus R_2$. Обов'язково зазначити, яке зображення матриці використовується в розв'язанні: двовимірний масив типу \textbf{ndarray} або список списків. Для всіх матриць має використовуватись однакове зображення.

21. Бінарні відношення $R_1$, $R_2$ задано на множині $\overline{0,n-1}$ матрицями відношення з елементами логічного типу. Обчислити матрицю відношення $R_1\cup R_2$. Обов'язково зазначити, яке зображення матриці використовується в розв'язанні: двовимірний масив типу \textbf{ndarray} або список списків. Для всіх матриць має використовуватись однакове зображення.

\newpage

22. Бінарні відношення $R_1$, $R_2$ задано на множині $\overline{0,n-1}$ матрицями відношення з елементами логічного типу. Обчислити матрицю відношення $R_1\cap R_2$. Обов'язково зазначити, яке зображення матриці використовується в розв'язанні: двовимірний масив типу \textbf{ndarray} або список списків. Для всіх матриць має використовуватись однакове зображення.

23. Бінарні відношення $R_1$, $R_2$ задано на множині $\overline{0,n-1}$ матрицями відношення з елементами логічного типу. Обчислити матрицю відношення $R_1\setminus R_2$. Обов'язково зазначити, яке зображення матриці використовується в розв'язанні: двовимірний масив типу \textbf{ndarray} або список списків. Для всіх матриць має використовуватись однакове зображення.

24. Бінарні відношення $R_1$, $R_2$ задано на множині $\overline{0,n-1}$ матрицями відношення з елементами логічного типу. Перевірити, чи правда, що  $R_1\subseteq R_2$. Обов'язково зазначити, яке зображення матриці використовується в розв'язанні: двовимірний масив типу \textbf{ndarray} або список списків.

25. Бінарні відношення $R_1$, $R_2$ задано на множині $\overline{0,n-1}$ матрицями відношення з елементами логічного типу. Перевірити, чи правда, що $R_1\subset R_2$. Обов'язково зазначити, яке зображення матриці використовується в розв'язанні: двовимірний масив типу \textbf{ndarray} або список списків.

26. Бінарні відношення $R_1$, $R_2$ задано на множині $\overline{0,n-1}$ матрицями відношення з елементами логічного типу. Обчислити композицію $R_1\circ R_2$ відношень. Обов'язково зазначити, яке зображення матриці використовується в розв'язанні: двовимірний масив типу \textbf{ndarray} або список списків.


27. Бінарне відношення $R$ задане на множині $\overline{0,n-1}$ матрицею відношення з елементами логічного типу. Перевірити, чи є відношення рефлексивним. Обов'язково зазначити, яке зображення матриці використовується в розв'язанні: двовимірний масив типу \textbf{ndarray} або список списків.

28. Бінарне відношення $R$ задане на множині $\overline{0,n-1}$ матрицею відношення з елементами логічного типу. Перевірити, чи є відношення антирефлексивним. Обов'язково зазначити, яке зображення матриці використовується в розв'язанні: двовимірний масив типу \textbf{ndarray} або список списків.

\newpage

29. Бінарне відношення $R$ задане на множині $\overline{0,n-1}$ матрицею відношення з елементами логічного типу. Перевірити, чи є відношення симетричним. Обов'язково зазначити, яке зображення матриці використовується в розв'язанні: двовимірний масив типу \textbf{ndarray} або список списків.

30. Бінарне відношення $R$ задане на множині $\overline{0,n-1}$ матрицею відношення з елементами логічного типу. Перевірити, чи є відношення антисиметричним. Обов'язково зазначити, яке зображення матриці використовується в розв'язанні: двовимірний масив типу \textbf{ndarray} або список списків.

31. Бінарне відношення $R$ задане на множині $\overline{0,n-1}$ матрицею відношення з елементами логічного типу. Перевірити, чи є відношення транзитивним. Обов'язково зазначити, яке зображення матриці використовується в розв'язанні: двовимірний масив типу \textbf{ndarray} або список списків.

32. Бінарне відношення $R$ задане на множині $\overline{0,n-1}$ матрицею відношення. Побудувати обернене відношення. Матрицю задано як список списків.

33. Бінарне відношення R задане на множині $\overline{0,n-1}$ матрицею відношення з цілими елементами.  Перевірити, чи є відношення рефлексивним. Обов'язково зазначити, яке зображення матриці використовується в розв'язанні: двовимірний масив типу \textbf{ndarray} або список списків.

34. Бінарне відношення R задане на множині $\overline{0,n-1}$ матрицею відношення з цілими елементами.  Перевірити, чи є відношення антирефлексивним. Обов'язково зазначити, яке зображення матриці використовується в розв'язанні: двовимірний масив типу \textbf{ndarray} або список списків.

35. Бінарне відношення R задане на множині $\overline{0,n-1}$ матрицею відношення з цілими елементами.  Перевірити, чи є відношення симетричним. Обов'язково зазначити, яке зображення матриці використовується в розв'язанні: двовимірний масив типу \textbf{ndarray} або список списків.

\newpage

36. Бінарне відношення R задане на множині $\overline{0,n-1}$ матрицею відношення з цілими елементами.  Перевірити, чи є відношення антисиметричним. Обов'язково зазначити, яке зображення матриці використовується в розв'язанні: двовимірний масив типу \textbf{ndarray} або список списків.

37. Бінарне відношення R задане на множині $\overline{0,n-1}$ матрицею відношення з цілими елементами.  Перевірити, чи є відношення транзитивним. Обов'язково зазначити, яке зображення матриці використовується в розв'язанні: двовимірний масив типу \textbf{ndarray} або список списків.

38. Написати функцію, що повертає кількість різних дійсних значень, що містяться в послідовності s.

39. За послідовністю цілих чисел s визначити яких значень в ній більше: парних чи непарних. Кожне значення враховується рівно один раз. Якщо більше парних, повернути 2, якщо більше непарних повернути 1, якщо кількості однакові, повернути 0. Розв’язок оформити у вигляді функції.

40. Словник a містить канонічний розклад натурального числа n на прості множники. Написати функцію, що за словником a обчислює число n.

41. Словники a1 та a2 містять канонічні розклади натуральних чисел n1 та n2 на прості множники. Написати функцію, що знаходить канонічний розклад найменшого спільного кратного цих чисел.

42. Словники a1 та a2 містять канонічні розклади натуральних чисел n1 та n2 на прості множники. Написати функцію, що знаходить канонічний розклад найбільшого спільного дільника цих чисел.

43. Словники a1 та a2 містять канонічні розклади натуральних чисел n1 та n2 на прості множники. Написати функцію, що знаходить найменше спільне кратне цих чисел.

44. Словники a1 та a2 містять канонічні розклади натуральних чисел n1 та n2 на прості множники. Написати функцію, що знаходить найбільший спільний дільник цих чисел.

45. Написати функцію, що за словником обчислює скільки разів зустрічається в ньому задане значення.

46. Написати функцію, що за послідовністю чисел будує список її елементів, що зустрічаються в ній рівно один раз. Числа, що мають різний тип, вважаємо різними.

\newpage

47. Написати функцію, що за послідовністю цілих знаходить 5 елементів, що зустрічаються в ній найчастіше. Якщо п’яте місце розділяють кілька значень, додавати до результату всіх їх. 

48. Написати функцію, що за послідовністю цілих знаходить елементи, що зустрічаються в ній рідше за інші (але зустрічаються). 

49. Написати функцію, що за послідовністю цілих знаходить всі елементи, що зустрічаються в ній рівно 7 разів. 

50. Написати функцію, що перевіряє чи складаються дві послідовності з однакових елементів з точністю до кількості входжень кожного елемента в послідовність. Порядок елементів у послідовності несуттєвий.

51. Список містить цілі числа. Вилучити зі списку усі повторення (кожне значення має залишитися в одному екземплярі). Для кожного числа має залишатись його перше входження в список.

52. Список містить числа. Вилучити зі списку усі повторення (кожне значення має залишитися в одному екземплярі). Для кожного числа має залишатись його перше входження в список. Числа, що мають різний тип, вважаємо різними.

53. Для функції $f$ типу $A\rightarrow B$ обернена функція, узагалі-то кажучи, функцією з $B\rightarrow A$ може не бути. Тим не менш можна розглядати "обернену" функцію $g$ з $B$ в $P(A)$ (булеан множини $A$), $g\left(b\right)=f^{-1}\left(b\right)$ ($f^{-1}\left(b\right)$ -- повний прообраз елемента $b$). Написати функцію, що за функцією $f$  типу $Z\rightarrow Z$, знаходить "обернену" до неї функцію $g$ як функцію з $Z$ в $P(Z)$. Вважати, що функцію $f$ задано словником, результатом функції також має бути словник.

54. Функція $f$ типу $A\times B \rightarrow C$ задана словником. Розглядати "обернену" функцію $g$ з $C \times B$  в $P(A)$ (булеан множини $A$), $g\left(c, b\right)=\left\{a\middle|\ f\left(a,b\right)=c\right\}$. Написати функцію, що за функцією $f$ типу $Z\times Z \rightarrow Z$, знаходить "обернену" до неї функцію $g$ як функцію з $Z\times Z$ в $P(Z)$. Вважати, що функцію $f$ задано словником, результатом функції також має бути словник.

55. За функцією $f$, заданою словником, та елементом $b$ знайти значення $f^{-1}(b)$. Розв’язання оформити у вигляді функції, що приймає $f$ та $b$ як аргументи.

\newpage

56. За функціями $f$ та $g$, заданими словниками знайти їх композицію. Розв’язання оформити у вигляді функції, що приймає словники $f$ та $g$ як аргументи та результат повертає у вигляді словника.

57. Для множин з повтореннями, заданими словниками, реалізувати теоретико-множинну операцію перевірки строгого включення. 

58. Для множин з повтореннями, заданими словниками, реалізувати теоретико-множинну операцію перевірки включення. Розв'яз\-ан\-ня оформити у вигляді функції.

59. Для множин з повтореннями, заданими словниками, реалізувати теоретико-множинну операцію об’єднання (елемент входить у результат максимум входжень в аргументи). Розв’язання оформити у вигляді функції.

60. Для множин з повтореннями, заданими словниками, реалізувати теоретико-множинну операцію
злиття (елемент входить в результат суму входжень в аргументи). Розв’язання оформити у вигляді функції.

61. Для множин з повтореннями, заданими словниками, реалізувати теоретико-множинну операцію перетину (елемент входить в результат мінімум входжень в аргументи). Розв’язання оформити у вигляді функції.

62. Для множин з повтореннями, заданими словниками, реалізувати теоретико-множинну операцію різниці. Розв’язання оформити у вигляді функції.

63. Відношення $C$ між множинами $A$ та $B$ ($A, B\subset Z$) задане словником, в якому ключами є елементи множини $A$, а значенням кожного ключа $a$ --- його повний образ $C\left(a\right)=\left\{b|\left(a,b\right)\in C\right\}$. Перевірити чи є відношення $C$ функціональним.

64. Відношення $C$ між множинами $A$ та $B$ ($A, B\subset Z$) задане словником, в якому ключами є елементи множини $A$, а значенням кожного ключа $a$ --- його повний образ $C\left(a\right)=\left\{b|\left(a,b\right)\in C\right\}$. Перевірити чи є відношення $C$ ін'єктивним.

65. Відношення $C$ між множинами $A$ та $B$ ($A, B\subset Z$) задане списком пар. Перевірити чи є відношення $C$ функціональним.

66. Відношення $C$ між множинами $A$ та $B$ ($A, B\subset Z$) задане списком пар. Перевірити чи є відношення $C$ ін'єктивним.

\newpage

67. Функція $f$ між множинами $A$ та $B$ (елементи множини $A$ є такими, що хешуються) задана списком пар. Написати функцію, що зображує $f$ у вигляді словника, де значеннями ключів є значення функції.

68. Відповідність $C$ між множинами $A$ та $B$ ($A\subset Z$) задана списком пар. Написати функцію, що зображує $C$ у вигляді словника, де значеннями ключів є їх повні образи.

69. Функція $f$ між множинами $A$ та $B$ (елементи множини $A$ є такими, що хешуються) задана словником, в якому значеннями ключів є значення функції. Написати функцію, що зображує $f$ у вигляді списку пар.

70. Бінарне відношення на множині $\overline{0,n-1}$ задане множиною пар. Перевірити, чи є відношення рефлексивним.

71. Бінарне відношення на множині $\overline{0,n-1}$ задане множиною пар. Перевірити, чи є відношення антирефлексивним.

72. Бінарне відношення на множині $\overline{0,n-1}$ задане множиною пар. Перевірити, чи є відношення симетричним.

73. Бінарне відношення на множині $\overline{0,n-1}$ задане множиною пар. Перевірити, чи є відношення антисиметричним.

74. Бінарне відношення на множині $\overline{0,n-1}$ задане множиною пар. Перевірити, чи є відношення транзитивним.

75. Бінарне відношення на множині $\overline{0,n-1}$ задане множиною пар. Побудувати обернене відношення.

76. Побудувати функцію, що зображення розбиття множини\\ $\overline{0,n-1}$  у вигляді множини множин перетворює на зображення розбиття у вигляді вектору (список або  \textbf{ndarray}), що задає належність елементів класам.

77. Побудувати функцію, що зображення розбиття множини\\ $\overline{0,n-1}$  у вигляді вектору (список або  \textbf{ndarray}), що задає належність елементів класам, перетворює на зображення у вигляді множини множин.

78. Побудувати функцію, що зображення розбиття множини у вигляді словника, що задає належність елементів класам, перетворює на зображення у вигляді множини множин.

\newpage

79. Побудувати функцію, що зображення розбиття множини у вигляді множини множин перетворює на зображення розбиття у вигляді словника, що задає належність елементів класам.

80. За еквівалентністю на множині $\overline{0,n-1}$ побудувати її фактор-множину як множину, елементами якої є множини.

81. За еквівалентністю на множині $\overline{0,n-1}$ побудувати її фактор-множину як вектор (список або  \textbf{ndarray}), що задає належність елементів класам.

82. За еквівалентністю на множині $A=\overline{0,n-1}$ та елементом множини $A$ обчислити його клас еквівалентності.

83. За розбиттям множини, що задане як множина множин, побудувати відповідну еквівалентність. Еквівалентність задати списком пар.

84. За розбиттям множини, що задане як множина множин, побудувати відповідну еквівалентність. Еквівалентність задати матрицею відношення.

85. За розбиттям множини $A=\overline{0,n-1}$, що задане як вектор (список або  \textbf{ndarray}), що задає належність елементів класам, побудувати відповідну еквівалентність. Еквівалентність задати списком пар.

86. За розбиттям множини $A=\overline{0,n-1}$, що задане як вектор (список або  \textbf{ndarray}), що задає належність елементів класам, побудувати відповідну еквівалентність. Еквівалентність задати матрицею відношення.

87. Визначити, чи є одне розбиття підрозбиттям іншого. Розбиття задані як множини множин.

88. Визначити, чи є одне розбиття підрозбиттям іншого. Розбиття задано як словники, що задають належність елементів класам.

89. Визначити, чи є одне розбиття множини $A=\overline{0,n-1}$\ підрозбиттям іншого. Розбиття задано векторами (список або  \textbf{ndarray}), що задають належність елементів класам.

\newpage

90. За розбиттям множини $A=\overline{0,n-1}$ та відображенням\\ $f{:}A\rightarrow{bool}$  побудувати підрозбиття таке, що два елементи знаходяться в одному класі підрозбиття тоді і тільки тоді, коли відображення $f$ приймає на них однакове значення. Вважати, що $f$ задано функцією. Для розбиттів використовувати зображення у вигляді множини множин.

91. За розбиттям множини $A=\overline{0,n-1}$ та відображенням $f$, заданим на множині $A$, побудувати підрозбиття таке, що два елементи знаходяться в одному класі підрозбиття тоді і тільки тоді, коли відображення $f$ приймає на них однакове значення. Вважати, що $f$ задано функцією. Розбиття зображувати векторами (список або  \textbf{ndarray}), що задають належність елементів класам.

92. Сімейство множин це множина, елементами якої є множини. Надалі будемо вважати, що сімейство множин задається як множина, елементами якої є незмінювані множини. Написати функцію, що за сімейством множин перевіряє, чи правда, що елементи сімейства попарно не перетинаються.

93. Сімейство множин це множина, елементами якої є множини. Надалі будемо вважати, що сімейство множин задається як множина, елементами якої є незмінювані множини. Написати функцію, що обчислює об'єднання сімейства множин. Об'єднання порожнього сімейства вважати порожньою множиною.

94. Сімейство множин це множина, елементами якої є множини. Надалі будемо вважати, що сімейство множин задається як множина, елементами якої є незмінювані множини. Написати функцію, що обчислює перетин сімейства множин. За спроби обчислити перетин порожнього сімейства має генеруватись виняток.

95. Сімейство множин це множина, елементами якої є множини. Надалі будемо вважати, що сімейство множин задається як множина, елементами якої є незмінювані множини. Написати функцію, що перевіряє, чи є задане сімейство множин розбиттям заданої множини $A$.

\newpage

96. Сімейство множин це множина, елементами якої є множини. Надалі будемо вважати, що сімейство множин задається як множина, елементами якої є незмінювані множини. Написати функцію, що перевіряє, чи є задане сімейство множин покриттям заданої множини $A$.




\end {Large}}
\end{document}
